%% sections/s3-amendment.tex - H. R. 9510 Bill v4.0 (full bill)
%% SEC. 3. AMENDMENT TO THE FEDERAL FOOD, DRUG, AND COSMETIC ACT.
%% Operative new section 515D (21 U.S.C. 360e-5), ten conforming amendments, the
%% clerical amendment, the rule of construction, and the effective date.
%% Gray-scale Mermaid figures (TikZ): Figure 3 (gate decision rule and funnel),
%% Figure 4 (statutory layering). Tables 2 and 3.
\clearpage
\billsec{SEC. 3.}{AMENDMENT TO THE FEDERAL FOOD, DRUG, AND COSMETIC ACT.}\phantomsection\label{toc:s3}

\uscprov{1}{\textbf{(a) In General.} Subchapter V of the Federal Food, Drug, and
Cosmetic Act (21 U.S.C. 351 et seq.) is amended by inserting after section 515C
(21 U.S.C. 360e-4) the following new section:}

\uscquote{0}{``SEC. 515D. VERIFICATION BEFORE GENERATION OF ROBOT-PATIENT
INTERACTION CODE IN PHYSICAL AI ONCOLOGY INVESTIGATIONS.}
\uscquote{0}{``(a) Requirement.- No robot-patient interaction code may be generated
or executed in a Physical AI oncology clinical investigation unless an automated
verification, validation, and uncertainty quantification process, bound to one or
more named external standards, has first cleared for that interaction, and the
cleared verification record has been documented and attested under subsection (e).}
\uscquote{0}{``(b) Order of Operations.- A sponsor shall not generate, and shall not
execute, robot-patient interaction code in advance of the cleared verification
record. A record created after generation or execution does not satisfy this
section.}
\uscquote{0}{``(c) Gate Thresholds and Decision Rule.- For each robot-patient
interaction, the verification process shall apply the gate thresholds set out in
the schedule in Table 2. Each gate shall emit a determination of ACCEPT, BLOCK, or
ESCALATE; a single failing dimension shall set BLOCK, and ambiguity or divergence
shall set ESCALATE and return the interaction to a qualified human. For each
interaction the verification fraction shall equal 1.0, and, for an interaction
subject to a hard catastrophe predicate, that predicate shall also hold
\cite{asmevv40, nasastd7009}.}

{\footnotesize
\begin{xltabular}{\textwidth}{@{}L{3.7cm} L{1.7cm} L{1.2cm} L{1.5cm} Y@{}}
\hline\noalign{\vskip 0.04cm}
\textbf{Gate} & \textbf{Min.\ agreement} & \textbf{Max.\ CV} & \textbf{Hard
predicate} & \textbf{Governing external standards}\\
\hline\noalign{\vskip 0.04cm}
\endfirsthead
\hline\noalign{\vskip 0.04cm}
\textbf{Gate} & \textbf{Min.\ agreement} & \textbf{Max.\ CV} & \textbf{Hard
predicate} & \textbf{Governing external standards}\\
\hline\noalign{\vskip 0.04cm}
\endhead
\hline\noalign{\vskip 0.04cm}
\endlastfoot
01 Bimanual hand-eye servo & 0.97 & 0.08 & No & IEC 80601-2-77; ISO 9283; ASME
V\&V 40\\
\hline\noalign{\vskip 0.04cm}
02 Dexterous finger force & 0.95 & 0.10 & No & IEC 80601-2-77; ISO/TS 15066\\
\hline\noalign{\vskip 0.04cm}
03 Whole-body balance & 0.98 & 0.06 & No & ISO 13482; IEC 60601-1\\
\hline\noalign{\vskip 0.04cm}
04 Auton. plan correct & 0.95 & 0.10 & No & UL 4600; IEEE 7009; IEC 62304\\
\hline\noalign{\vskip 0.04cm}
05 Instr. grasp, handover & 0.96 & 0.10 & No & IEC 80601-2-77; ISO 9283\\
\hline\noalign{\vskip 0.04cm}
06 Vasc. no-fly vol. (hand) & 1.00 & 0.05 & Yes & IEC 80601-2-77; ISO 14971\\
\hline\noalign{\vskip 0.04cm}
07 Bimanual suturing anast. & 0.96 & 0.08 & No & IEC 80601-2-77; Fistula Risk
Score\\
\hline\noalign{\vskip 0.04cm}
08 Perception scene unders. & 0.95 & 0.10 & No & ASME V\&V 40; IEC 62304\\
\hline\noalign{\vskip 0.04cm}
09 Shared-room human coll. & 1.00 & 0.06 & Yes & ISO/TS 15066; ISO 10218-1; ISO
13482\\
\hline\noalign{\vskip 0.04cm}
10 Fault, emer. stop, degrad. & 1.00 & 0.05 & Yes & IEC 60601-1; ISO 13849-1; IEEE
7009\\
\hline
\end{xltabular}}
\vspace{-0.8cm}
\figcaption{Table 2. Ten-gate threshold schedule: min. validation agreement, max.
coefficient of variation, hard catastrophe predicate, and external standards for
each gate. Gates 06, 09, and 10 are hard predicate at agreement 1.00.}

\uscprov{1}{The decision rule and the gate funnel that this schedule implements are
shown in Figure 3.}

\begin{mermaidfig}
\node[mmstep,text width=40mm] (g) at (0,0) {Candidate robot-patient interaction code};
\node[mmdec,text width=24mm] (q) at (0,-2.0) {Clears all 10 gates?};
\node[mmgoal,text width=26mm] (a) at (3.8,-2.0) {ACCEPT: code may be generated};
\node[mmdec,text width=20mm] (e) at (0,-4.2) {Borderline?};
\node[mmlaw,text width=24mm] (b) at (3.8,-4.2) {BLOCK: code withheld};
\node[mmstep,text width=26mm] (esc) at (-3.8,-4.2) {ESCALATE to human review};
\draw[mmedge] (g) -- (q);
\draw[mmedge] (q) -- node[mmlabel]{yes} (a);
\draw[mmedge] (q) -- node[mmlabel]{no} (e);
\draw[mmedge] (e) -- node[mmlabel]{no} (b);
\draw[mmedge] (e) -- node[mmlabel]{yes} (esc);
\draw[mmedge] (esc) |- (g);
\end{mermaidfig}
\vspace{-0.3cm}
\figcaption{Figure 3. The VVUQ gate decision rule over the verify, validate, and
quantify dimensions: a single failing dimension sets BLOCK, divergence sets
ESCALATE to a qualified human, and only a full clear sets ACCEPT.}

\uscquote{0}{``(d) Readiness Gates.- Before a patient is enrolled or a robot acts,
the trial site shall pass the Physical AI Standard Level gate on legislative
authorization, regulatory compliance, and operational readiness; each robot shall
meet the minimum Unification Standard Level for its functional type, which is not
less than 7.0 for a surgical type, 6.0 for a therapeutic positioning or diagnostic
needle placement type, 4.0 for a rehabilitative type, and 3.0 for a companion
monitoring type; and each robot's behavior shall be validated across not fewer than
two simulation frameworks within the cross-framework tolerances of less than 2
millimeters of trajectory deviation and less than 0.5 newton of force discrepancy
\cite{kawchak2026national}.}
\uscquote{0}{``(e) Documentation and Attestation.- Before robot-patient interaction
code is generated or executed, the sponsor shall make and file, for each covered
algorithm, algorithm documentation including a plain-language summary, the decision
logic or architecture, the training or conditioning data sets, the gate bindings
identifying the governing external standard and the threshold enforced, the
verification-before-generation record, and a determinism statement giving the seed;
a compliance statement and a signed attestation by a responsible officer; and a
hash-chained, tamper-evident audit trail that satisfies part 11 of title 21, Code
of Federal Regulations \cite{cfr-part11, cfr-hipaa-164}.}
\uscquote{0}{``(f) Cybersecurity, Human Oversight, and Lifecycle.- Each Physical AI
system shall incorporate authentication and access control, encryption of data in
transit and at rest, network segmentation, and software-integrity verification;
shall operate under recorded human oversight with a manual override and an
emergency stop that brings the system to a safe state within the
procedure-appropriate budget and a hand-back-to-human escalation when a gate
escalates or a fault is detected; and shall maintain configuration and
model-version management, monitor for model and data drift, and, on
decommissioning, archive the audit trail, securely erase patient data, and revoke
access \cite{iec60601, iso13849}.}
\uscquote{0}{``(g) Nondiscrimination.- Each covered algorithm and its training data
shall minimize the risk of bias on protected characteristics and adhere to
evidence-based clinical guidelines, and each Physical AI system shall be designed
and operated consistent with applicable accessibility and accommodation
requirements \cite{cfr-92-210, usc-titlevi, usc-ada}.}
\uscquote{0}{``(h) Regulations.- Not later than 365 days after the date of enactment
of this section, the Secretary shall issue final regulations to carry out this
section. Pending those regulations, the Secretary may issue interim guidance, and a
sponsor may comply by meeting the requirements of subsections (a) through (g).}
\uscquote{0}{``(i) Definitions.- In this section, the terms `Physical AI system',
`robot agent', `robot-patient interaction code', `generative artificial
intelligence', `AI-enabled device software function', `verification', `validation',
`uncertainty quantification', `verification fraction', `validation agreement',
`coefficient-of-variation bound', `hard catastrophe predicate', `Physical AI
Standard Level', `Unification Standard Level', and `hash-chained audit trail' have
the meanings given in the prior bill and reconciled with section 201(h) and section
520(o) of this current Act.}
\uscquote{0}{``(j) Rule of Construction.- Nothing in this section displaces the
investigational-device exemption under section 520(g), the protection of human
subjects under part 50 of title 21, Code of Federal Regulations, the
investigational new drug framework under part 312 of such title, or the
predetermined change control plan authority under section 515C as it applies
outside Physical AI, and nothing in this section itself reclassifies any device.''.}

\uscprov{1}{\textbf{(b) Conforming Amendments.} The Federal Food, Drug, and Cosmetic
Act is amended, with the corresponding changes shown in the comparative print in
section 4, so that the device definition in section 201(h) reaches software that
generates or executes robot-patient interaction code; the prohibited acts in
section 301 reach a violation of section 515D; a device used without a cleared
section 515D record is adulterated under section 501; the real world evidence
program in section 505F recognizes section 515D records; a verification-governed
change consistent with a predetermined change control plan needs no new premarket
notification under section 510(k); special controls under section 513 are graduated
by degree of autonomy; a premarket approval application under section 515 must
include the cleared verification record; a predetermined change control plan under
section 515C must include an automated verification-before-change protocol; the
clinical-decision-support exclusion in section 520(o) does not reach robot-control
software; and a savings clause at section 521 preserves State human-in-the-loop
requirements that are not different from, or in addition to, a Federal device
requirement \cite{usc-321h, usc-331, usc-351, usc-355g, usc-360, usc-360c,
usc-360e, usc-360e4, usc-360j, usc-360k}. Figure 4 shows where the new section
attaches on the device-regulation pathway, and Table 3 is the conforming-amendment
crosswalk.}

\begin{mermaidfig}
\node[mmin,text width=30mm]  (d)  at (0,0)    {Is it a device? sec. 321(h)};
\node[mmin,text width=34mm]  (cds)at (4.6,0)  {CDS exclusion sec. 360j(o) does NOT carve out robot-control software};
\node[mmstep,text width=30mm](c)  at (0,-1.7) {Classify the device sec. 360c};
\node[mmstep,text width=26mm](k)  at (-2.8,-3.4){Class II 510(k) sec. 360(k)};
\node[mmstep,text width=26mm](p)  at (2.8,-3.4) {Class III PMA sec. 360e};
\node[mmlaw,text width=34mm] (pc) at (0,-5.1) {PCCP keystone sec. 360e-4};
\node[mmgoal,text width=44mm](n)  at (0,-6.8) {NEW sec. 360e-5: mandatory, ex ante, automated VVUQ rule};
\draw[mmedge] (d) -- (c);
\draw[mmedged] (cds) -- (d);
\draw[mmedge] (c) -- (k);
\draw[mmedge] (c) -- (p);
\draw[mmedge] (k) -- (pc);
\draw[mmedge] (p) -- (pc);
\draw[mmedge] (pc) -- (n);
\end{mermaidfig}
\vspace{-0.3cm}
\figcaption{Figure 4. Statutory layering of new section 360e-5 through Title 21: the
device question, the clinical-decision-support carve-out it must not fall through,
classification, the 510(k) and premarket-approval pathways, the section 360e-4
keystone, and the new mandatory, ex ante section 360e-5.}

\clearpage

{\footnotesize
\begin{xltabular}{\textwidth}{@{}L{1.0cm} L{1.6cm} L{1.7cm} Y@{}}
\hline\noalign{\vskip 0.04cm}
\textbf{No.} & \textbf{FD\&C \S} & \textbf{21 U.S.C.} & \textbf{Conforming change
made by H. R. 9510}\\
\hline\noalign{\vskip 0.04cm}
\endfirsthead
\hline\noalign{\vskip 0.04cm}
\textbf{No.} & \textbf{FD\&C \S} & \textbf{21 U.S.C.} & \textbf{Conforming change
made by H. R. 9510}\\
\hline\noalign{\vskip 0.04cm}
\endhead
\hline\noalign{\vskip 0.04cm}
\endlastfoot
(1) & \S~201(h) & \S~321(h) & Device definition reaches software that generates or
executes robot-patient code\\
\hline\noalign{\vskip 0.04cm}
(2) & \S~301 & \S~331 & New prohibited act for generation or execution in violation
of section 515D\\
\hline\noalign{\vskip 0.04cm}
(3) & \S~501 & \S~351 & Device used without a cleared section 515D record is
adulterated\\
\hline\noalign{\vskip 0.04cm}
(4) & \S~505F & \S~355g & Real world evidence program recognizes section 515D
records\\
\hline\noalign{\vskip 0.04cm}
(5) & \S~510(k) & \S~360(k) & Verification-governed change needs no new premarket
notification\\
\hline\noalign{\vskip 0.04cm}
(6) & \S~513 & \S~360c & Special controls graduated by degree of autonomy\\
\hline\noalign{\vskip 0.04cm}
(7) & \S~515 & \S~360e & Premarket approval application must include the cleared
verification record\\
\hline\noalign{\vskip 0.04cm}
(8) & \S~515C & \S~360e-4 & Change control plan must include automated
verification-before-change protocol\\
\hline\noalign{\vskip 0.04cm}
(9) & \S~520(o) & \S~360j(o) & Clinical-decision-support exclusion does not reach
robot-control software\\
\hline\noalign{\vskip 0.04cm}
(10) & \S~521 & \S~360k & Savings clause preserves State human-in-the-loop
requirements\\
\hline
\end{xltabular}}
\vspace{-0.8cm}
\figcaption{Table 3. The ten conforming amendments crosswalk: each amendment by its
Federal Food, Drug, and Cosmetic Act section, its 21 U.S.C. citation, and the
change H. R. 9510 makes. The text appears in comparative print section 4.}

\uscprov{1}{\textbf{(c) Clerical Amendment.} The table of contents for the Federal
Food, Drug, and Cosmetic Act is amended by inserting after the item relating to
section 515C the following new item:}
\uscquote{0}{``Sec. 515D. Verification before generation of robot-patient
interaction code in Physical AI oncology investigations.''.}

\uscprov{1}{\textbf{(d) Rule of Construction.} Nothing in this Act displaces the
investigational-device exemption under section 520(g) of the Federal Food, Drug, and
Cosmetic Act (21 U.S.C. 360j(g)), the protection of human subjects under part 50 of
title 21, Code of Federal Regulations, the investigational new drug framework under
part 312 of such title, the predetermined change control plan authority under
section 515C as it applies outside Physical AI, or any State law preserved under
section 521 as amended; and nothing in this Act reclassifies any device
\cite{cfr-part50, cfr-part312}.}

\uscprov{1}{\textbf{(e) Effective Date; Implementation Deadline; Transition Rule.}
This Act takes effect on the first day of the first calendar quarter beginning not
less than one year after enactment; the Secretary shall issue the final regulations
required by section 515D(h) not later than 365 days after enactment; and a Physical
AI oncology clinical investigation already enrolling on the effective date shall
have 180 days to come into compliance with section 515D, and this Act does not
disturb the February 2, 2026, effective date of the Quality Management System
Regulation \cite{fr-qmsr-2024}.}
